\documentclass[11pt,letterpaper]{article}
\usepackage[margin=0.9in,headheight=14pt]{geometry}
\usepackage[T1]{fontenc}
\usepackage{lmodern}
\usepackage{amsmath,amssymb,mathtools}
\usepackage{microtype,xcolor,booktabs,array,enumitem,fancyhdr}
\usepackage[colorlinks=true,linkcolor=navy,citecolor=navy,urlcolor=navy]{hyperref}
\definecolor{navy}{RGB}{27,52,80}
\definecolor{teal}{RGB}{20,92,89}
\definecolor{wine}{RGB}{128,39,56}
\hypersetup{pdftitle={Integral images of Springer--Borel maps: relation to Parts II and III},
pdfauthor={Discussion note prepared for Shrawan Kumar},
pdfsubject={Audited general type B Smith formula, published inputs, and remaining type D questions}}
\pagestyle{fancy}
\fancyhf{}
\fancyhead[L]{\small Integral Springer--Borel images}
\fancyhead[R]{\small Discussion note}
\fancyfoot[C]{\small\thepage}
\renewcommand{\headrulewidth}{0.3pt}
\setlength{\parindent}{0pt}
\setlength{\parskip}{5pt plus 1pt}
\setlist[itemize]{leftmargin=1.5em,itemsep=3pt,topsep=3pt}
\setlist[enumerate]{leftmargin=1.7em,itemsep=3pt,topsep=3pt}
\setlength{\emergencystretch}{2em}
\newcommand{\Z}{\mathbb Z}
\newcommand{\Q}{\mathbb Q}
\newcommand{\C}{\mathbb C}
\newcommand{\F}{\mathbb F}
\newcommand{\OG}{\operatorname{OG}}
\newcommand{\IG}{\operatorname{IG}}
\newcommand{\Gr}{\operatorname{Gr}}
\newcommand{\sat}{\operatorname{sat}}
\newcommand{\rank}{\operatorname{rank}}
\newcommand{\Smith}{\operatorname{Smith}}
\newcommand{\qbinom}[2]{\genfrac{[}{]}{0pt}{}{#1}{#2}_q}
\newcommand{\published}{\textcolor{navy}{\textbf{Published result/input.}}\ }
\newcommand{\deduction}{\textcolor{teal}{\textbf{Deduction in the companion notes.}}\ }
\newcommand{\conjectural}{\textcolor{wine}{\textbf{Conjecture.}}\ }
\title{\vspace{-1.5em}\textcolor{navy}{Integral images of Springer--Borel maps}\\[4pt]
\large Connection with Parts II and III}
\author{Discussion note prepared for Shrawan Kumar}
\date{19 October 2026}

\begin{document}
\maketitle
\thispagestyle{empty}
\vspace{-1.0em}

\textbf{Purpose and status.}
This note summarizes an integral refinement of the maps studied in
Kumar--Rogers, Part II \cite{KR}, and Kumar--Xie, Part III \cite{KX}.
It separates published statements, deductions whose arguments are recorded
and checked in the companion repository, and unproved conjectures.
The general normalized type-B formula now has a complete argument in the
companion notes, freshly audited here and written out below.
No gap was found in that review, including the earlier integral reduction.
This is an unpublished repository deduction; external expert review and
originality remain unestablished.

\section{The published starting point}
\published Part II gives the classical symplectic and odd-orthogonal
generator images and stable injectivity: Theorem 16 for C and Theorem 29
for B. Part III gives the even-orthogonal generator images, its stable
integral presentation, and stable injectivity, especially Theorems 3.6 and
3.8 and Proposition 3.7. The injectivity theorems are stated over \(\C\);
their proofs explicitly use chosen integral forms.

Use the normalized torus coordinates
\(\bar t_j=(t_j-t_j^{-1})/2\), corresponding to \(x_j\) under the
Springer--Borel map. Put \(c_i=c_i(Q)\), where \(Q\) is the tautological
quotient bundle on the relevant isotropic Grassmannian, with \(c_0=1\)
and \(c_i=0\) outside its rank. The squared-tail generator image is
\[
 e_i(\bar t_{\rm tail}^{\,2})
 \longmapsto c_i^2+
 2\sum_{j=1}^{i}(-1)^j c_{i+j}c_{i-j}.
 \tag{1}
\]
The integral theta/eta presentations and B/C Schubert comparison used
below are existing Schubert calculus \cite{BKTIG,BKTD}.

\section{The additional integral question}
Write \(n=a+k\). On
\(X_B=\OG(a,2n+1)\) use
\(q_i=e_i(x_{a+1}^2,\ldots,x_n^2)\), \(1\le i\le k\).
On \(X_D=\OG(a,2n)\) use \(q_1,\ldots,q_{k-1}\) and
\(p=x_{a+1}\cdots x_n\), with \(p^2=q_k\).
Take \(a,k\ge1\) in B and \(a\ge1,k\ge2\) in D.

Polynomial degree \(d\) means cohomological degree \(2d\);
\(\deg q_i=2i\) and \(\deg p=k\).
Let \(L_d\subset H^{2d}(X;\Z)\) be the integer span of all degree-\(d\)
monomials in these fixed inputs after applying the Borel map. The object
to determine is
\[
 \sat(L_d)/L_d,\qquad
 \sat(L_d)=(L_d\otimes\Q)\cap H^{2d}(X;\Z).
 \tag{2}
\]
We retain every nonzero Smith factor of \(L_d\hookrightarrow\sat(L_d)\),
including \(1\): this means the finite group (2) together with \(\rank L_d\).
The index is in the selected rational image, whose rank can be smaller
than ambient cohomology. Injectivity alone does not imply saturation.


\newpage
\section{The general normalized type-B result}
Put
\[
 \ell=\min(a,k),\quad A=\max(a,k),\quad
 r=\lfloor n/2\rfloor,\quad c=\lceil n/2\rceil,\quad
 \epsilon=c-r\in\{0,1\}.
\]
Let \(M_e(m)\) count nonzero factors \(2^e\) in polynomial degree \(2m\),
equivalently cohomological degree \(4m\).
The variable \(q\) counts \(m\), not the generator \(q_i\).
Write \(G(v,u)=\qbinom{v}{u}\), zero outside \(0\le u\le v\).
Define \(U_{0,0}=1\) and
\[
 U_{n,L}(q)=
 \sum_{j=0}^{\lfloor L/2\rfloor}
 q^{j(c-L+j)}G(r,j)G(c-j,L-2j).
 \tag{3}
\]

\deduction For the fixed normalized inputs of Section 2,
\[
 \boxed{\quad
 \sum_m M_e(m)q^m
 =q^{ec}G(r,e)\,U_{n-2e,\ell-e}(q),
 \qquad 0\le e\le\ell .
 \quad} \tag{4}
\]
Recompute \(r,c\) inside \(U\) at its displayed arguments.
There is no image in odd polynomial degree.
All nonzero factors are powers of two, bounded by \(2^\ell\).
The formula determines their multiplicities in every rank and degree.

Here is an equivalent partition rule. For
\(\gamma\subseteq A^\ell\), pad to \(\ell\) parts and put
\[
 \beta_i=\gamma_i+\ell-i,\qquad
 d(\gamma)=\#\{i:\beta_i<r\},\qquad
 e(\gamma)=\#\{i:\beta_i\ge c+d(\gamma)\}.
 \tag{5}
\]
In degree \(2m\), the Smith multiset is
\(\{2^{e(\gamma)}:|\gamma|=m,\ \gamma\subseteq A^\ell\}\).
This labels multiplicities; it does not assert that the original Schur
generators are individually Smith vectors.

\paragraph{Examples and endpoints.}
For \(a=k=2\) in polynomial degree \(4\), the factors are \((1,2)\).
For \(a=k=4\), the counts over all degrees are \(19,28,18,4,1\)
for factors \(1,2,4,8,16\), totaling \(\binom84=70\).
The two endpoints of (4) are
\[
 \sum_m M_0(m)q^m=U_{n,\ell}(q),\qquad
 \sum_m M_\ell(m)q^m=q^{\ell c}G(r,\ell).
 \tag{6}
\]

\paragraph{Connection to the papers.}
Parts II and III determine the maps and prove stable complex injectivity.
Formula (4) addresses additional finite integral information: the index
of the selected image in its saturation, degree by degree.
The companion notes also deduce stable normalized saturation in B and D,
and sufficient finite saturation ranges \(d\le A\) in B and \(d\le a\)
in D. These integral conclusions are not attributed to the papers as
published Smith theorems.

\newpage
\section{The integral reduction to Schur prefixes}
We give the argument for (4). The full source and dependency audit are
linked in \cite{Repo}. By integral block duality put the smaller block
at the head. In polynomial degree \(2m\), use the source basis
\(s_\lambda(y_1^2,\ldots,y_\ell^2)\), with
\(\lambda\subseteq A^\ell\), \(|\lambda|=m\).
Inversion of the two squared elementary series and Schur transposition
give this integral change from the tail generators.

\published The Chern-preserving B/C Schubert comparison
\cite{BKTIG} gives coefficient matrices
\[
 B_{w,\lambda}=2^{\nu(w)}C_{w,\lambda},\qquad
 \nu(w)=\#\{\text{negative signed entries}\}.
 \tag{7}
\]

\deduction The selected C image is primitive, by the integral quaternionic Schur
basis and the isotropic-to-quaternionic fibration \cite{PR}.
Thus its full row lattice in these source coordinates is \(\Z^R\).

Identify that coefficient lattice with degree \(m\) in
\[
 R_{A,\ell}=\Z[e_1,\ldots,e_\ell]/
             (h_{A+1},\ldots,h_{A+\ell})
\]
by the rectangle Schur basis, and write
\(P_t=\operatorname{span}_{\Z}\{C_{w,\bullet}:\nu(w)\le t\}\);
\(P_{-1}=0\), \(P_\ell=R_{A,\ell}\).
This identifies coefficient rows, not cohomological degrees of the
original characteristic map.

\deduction The integral prefix reduction is
\[
 P_t=\operatorname{span}_{\Z}\{s_\mu s_\alpha s_\beta\},\quad
 \begin{cases}
 \mu\subseteq A^t,\\
 \alpha\subseteq(c-p)^p,\\
 \beta\subseteq(r-s)^s,
 \end{cases}
 \quad
 b=\ell-t,\ p=\lceil b/2\rceil,\ s=\lfloor b/2\rfloor .
 \tag{8}
\]
Products are in \(R_{A,\ell}\), of the degree under consideration.

Here are the integral points in that reduction. The cells with
\(\nu\le t\) form the closed locus
\(\dim(\Sigma\cap L_0)\ge b\) in \(\IG(\ell,2n)\).
Choosing a \(b\)-plane in this intersection resolves the locus.
Root-group coordinates supply a section over every Schubert cell;
the graph closure lifts its closure with degree one.
Hence cohomology pullback is primitive over \(\Z\).

The resolution maps to a topological quaternionic Grassmann bundle
\(Q\to\Gr(b,n)\); the map to \(Q\) has Lagrangian fiber.
Global quaternionic characteristic classes give an integral
Leray--Hirsch basis for \(Q\) over \(\Gr(b,n)\) \cite{PR}.
The algebraic relative analogue is \cite{PW}.
The even integral Serre spectral sequence makes pullback from \(Q\)
primitive.
Branching and Littlewood's two-quotient identity \cite{Wildon}
give (8): empty 2-core beta numbers in \(\{0,\ldots,n-1\}\)
have \(p,s\) beads on the two runners, giving exactly those rectangles.
Both primitive pullbacks have integral left inverses, preserving entire
row lattices. The fiber width is \(A\); no hidden Schur reduction occurs.

If all \(P_t\) are primitive, an adapted integral basis gives
multiplicity \(\rank P_t-\rank P_{t-1}\) for \(2^t\), since the
weighted row lattice in (7) is \(\sum_t2^tP_t\).
The rest of the proof establishes precisely that hypothesis and its ranks.

\newpage
\section{A geometric rank bound}
Let \(N:V\to V\), \(V=\C^n\), be square zero of maximal rank \(r\).
Put \(I=\operatorname{im}N\), \(K=\ker N\), of dimensions \(r,c\), and
\[
 W_t=\{U\in\Gr(\ell,V):
              \delta(U):=\rank(N:U\to V/U)\le t\}.
\]
This is a projective determinantal locus.
For the stratum of defect \(e\), put
\(D=NU\), \(B=U\cap K\), \(C=D\cap B\), \(j=\dim C\).
Then \(\dim D=e+j\) and \(\dim B=\ell-e-j\).
Choose successively
\[
 \begin{gathered}
 D\in\Gr(e+j,r),\qquad C\in\Gr(j,D),\\
 (B/C)\bmod D\in\Gr(\ell-e-2j,K/D).
 \end{gathered}
\]
The two graph lifts have dimensions \(e(\ell-e-2j)\) and
\((e+j)(c-\ell+e+j)\). Their sum is \(ec+j(c-\ell+j)\).
Therefore the compact-support even Betti polynomial of this stratum is
\[
 \begin{aligned}
 E_e^{\rm geom}(q)
 &=\sum_j q^{ec+j(c-\ell+j)}
       G(r,e+j)G(e+j,j)G(c-e-j,\ell-e-2j)\\
 &=q^{ec}G(r,e)\,U_{n-2e,\ell-e}(q).
 \end{aligned} \tag{9}
\]
Terms with impossible dimensions are zero.
Filter first by defect, then by \(\rank(N|_U)=e+j\);
these are closed rank loci in the successive strata.
Grassmann cell refinements and the affine lifts give free even integral
homology, generated by algebraic cycle closures.
The projectivity of \(W_t\) gives
\[
 \sum_m\rank H^{2m}(W_t;\Z)q^m=\sum_{e=0}^t E_e^{\rm geom}(q).
 \tag{10}
\]

For \(b,p,s\) as in (8), take the smooth projective iterated
Grassmann bundle
\[
 Y_{\rm bal}=\{(F,S,U):
 F\subset I,\ \dim F=s;\
 F\subset S\subset N^{-1}F,\ \dim S=b;\
 S\subset U,\ \dim U=\ell\}.
 \tag{11}
\]
Over \(\Gr(s,r)\), \(E=N^{-1}F/F\) has rank \(c\) and
\(c(E)=1\), since \(0\to K/F\to E\to F\to0\).
The next two fibers are \(\Gr(p,E)\) and
\(\Gr(t,V/S)\), of widths \(c-p\) and \(A\).
Schur branching gives exactly the rows (8).
Since \(S\) is invariant, the forgetful map factors through \(W_t\).
Consequently
\[
 \rank P_t(m)\le\rank H^{2m}(W_t;\Z).
 \tag{12}
\]
This uses factorization of pullback, without assuming ambient Chern
generation or surjectivity of restriction to \(W_t\).

\newpage
\section{The integral bridge: a Chow envelope}
The proper map \(f:Y_{\rm bal}\to W_t\) has degree-one lifts of
every irreducible subvariety, not just lifts of complex points.

Indeed the following construction works over every extension field of
\(\C\). For \(U\in W_t\), the invariant space
\(U\cap N^{-1}U\) has dimension at least \(b\).
A square-zero Jordan basis chooses an invariant \(b\)-plane \(S\)
inside it, over the same field.
Writing \(j=\rank N|_S\le s\), we have
\[
 \dim(S\cap I)\ge\max(j,b-j-\epsilon)\ge s.
\]
For the last inequality use the first term if \(j=s\);
otherwise \(b-j-\epsilon\ge b-s\ge s\).
As \(NS\subset S\cap I\), choose \(NS\subset F\subset S\cap I\)
of dimension \(s\). This supplies a point of (11) above \(U\).

Apply this at the function field of an irreducible
\(Z\subset W_t\). Closing the graph of its rational lift gives a
subvariety of \(Y_{\rm bal}\) birational onto \(Z\).
Thus \(f\) is a Chow envelope:
\[
 f_*:\operatorname{CH}_*(Y_{\rm bal})
       \twoheadrightarrow\operatorname{CH}_*(W_t)
 \quad\hbox{over }\Z .
 \tag{13}
\]
The same cycle-lifting and the free even homology in Section 5 also make
\(f^*\) a primitive cohomology injection. Only (13) is needed for
the following integral containment.

\paragraph{Unbalanced incidence rows.}
For every \(0\le j\le b\), define \(Y_j\) by (11) with
\(\dim F=b-j\) and \(\dim S/F=j\).
All dimensions are permitted because \(b\le\ell\le r\).
The map \(g_j:Y_j\to W_t\) has coefficient rows
\[
 s_\mu s_\alpha s_\beta,\qquad
 \mu\subseteq A^t,\quad
 \alpha\subseteq(c-j)^j,\quad
 \beta\subseteq(r-b+j)^{b-j}.
 \tag{14}
\]
Out-of-width middle classes vanish because \(c(E)=1\).
The top factors already fit width \(A\), being subpartitions of the
source partition. Thus branching introduces no omitted bundle terms.

Every row in (14) belongs to \(P_t\) over \(\Z\).
In degree \(2m\), choose an integral cycle
\(D\in\operatorname{CH}_m(Y_j)\) dual to the relevant cohomology basis
element. Such cycles exist by the integral Grassmann-bundle formula
and perfect intersection pairing on the smooth \(Y_j\).
Lift \((g_j)_*D\) through (13). The proper projection formula expresses
its pairings with ambient Schur classes as an integer combination of
the balanced rows (8).
The target \(W_t\) may be singular: the argument uses Chow cycles and
ambient Chern classes acting on them, so needs no cohomological
pushforward on \(W_t\).
This is the step that upgrades the calculation from rational ranks to
an integral lattice comparison.

\newpage
\section{Primitive pivots and the bead criterion}
The complete-symmetric products \(h_\gamma\),
\(\gamma\subseteq A^\ell\), form an integral basis of \(R_{A,\ell}\).
Kostka/Jacobi--Trudi triangularity gives this change of basis;
\(h_i=0\) for \(i>A\).
The same bounded change lets us replace \(s_\mu\) by \(h_\mu\)
in (14).

For \(0\le j\le\lfloor(b+\epsilon)/2\rfloor\), put
\[
 \Delta_j=((c-j)^j,(r-b+j)^{b-j}),\qquad
 \mathcal K_t=\bigcup_j\{\gamma:\gamma\subseteq(A^t,\Delta_j)\}.
 \tag{15}
\]
The corners are weakly decreasing and fit \(A^\ell\).
For \(b=0\), take \(\mathcal K_\ell=A^\ell\).
Given a pivot \(\gamma\), choose a corner, split its parts into
blocks \(\mu,\alpha,\beta\) of lengths \(t,j,b-j\), and use the
integral row \(v_\gamma=h_\mu s_\alpha s_\beta\) from (14).

Its coefficient on \(h_\gamma\) is one.
Every other surviving complete coordinate has strictly larger energy
\(\sum_i\gamma_i^2\). For example, a nonidentity Jacobi--Trudi
permutation changes the \(\alpha\) block's energy by
\[
 2\sum_i\alpha_i(\sigma(i)-i)
       +\sum_i(\sigma(i)-i)^2>0.
 \tag{16}
\]
The first sum is nonnegative by rearrangement and the second positive.
Energies add between the two blocks; invalid indices vanish.
Order by energy and project these rows onto the columns \(\mathcal K_t\).
The square minor is triangular with diagonal one.
Their span \(Q_t(m)\subseteq P_t(m)\) is therefore primitive, of rank
\(\#\{\gamma\in\mathcal K_t:|\gamma|=m\}\).

\paragraph{Exact membership.}
For the beta numbers and statistic (5),
\[
 \gamma\in\mathcal K_t\quad\Longleftrightarrow\quad e(\gamma)\le t.
 \tag{17}
\]
Here is the boundary-sensitive argument. If \(b=0\), it is automatic.
Otherwise put \(d=\#\{\beta_i<r\}\) and \(J=\ell-d\).
If \(J\le t\), all bottom \(b\) beads are below \(r\);
the \(j=0\) corner works, and \(e(\gamma)\le J\le t\).
If \(J>t\), those bottom beads have \(j_0=J-t=b-d>0\) entries
at least \(r\). The nonempty corner bounds are
\[
 J\le t+j,\qquad \beta_{t+1}\le c+b-j-1.
\]
Any permissible corner has \(j\ge j_0\), hence
\(\beta_{t+1}<c+d\), equivalent to \(e(\gamma)\le t\).
Conversely that inequality puts its \(j_0\) high beads in
\([r,c+d-1]\), of length \(d+\epsilon=b-j_0+\epsilon\).
Thus \(2j_0\le b+\epsilon\), and the corner \(j_0\) works.
This includes an empty middle interval and \(b=1\).

\newpage
\section{Weighted counting and the integral conclusion}
For partitions of exact statistic \(e\), choose \(d\) beads below \(r\),
\(j=\ell-d-e\) beads in \([r,c+d-1]\), and \(e\) above that interval.
Weighted subset counting contributes
\(G(r,d)G(d+\epsilon,j)G(r-d,e)\).
Since \(|\gamma|=\sum_i\beta_i-\ell(\ell-1)/2\), its minimum-degree
shift is \(ec+j(r-\ell+j)\).
Set \(R=r-e\), \(L=\ell-e\); the flag identity yields
\[
 \sum_{\gamma:e(\gamma)=e}q^{|\gamma|}
 =q^{ec}G(r,e)
   \sum_{j\ge0}q^{j(R-L+j)}
             G(R,L-j)G(L-j+\epsilon,j).
 \tag{18}
\]
For \(\epsilon=0\), the inner product becomes
\(G(R,j)G(R-j,L-2j)\), so the inner sum is \(U_{2R,L}\).
For \(\epsilon=1\), use
\[
 G(L-j+1,j)=G(L-j,j-1)+q^jG(L-j,j).
\]
Shift \(j=s+1\) in the first term, apply the same flag identity,
and combine with the second using
\[
 G(R+1-s,L-2s)=G(R-s,L-2s)
       +q^{R+1-L+s}G(R-s,L-2s-1).
\]
This gives \(U_{2R+1,L}\). Thus (18) is exactly (4).

Summing over \(e\le t\), the pivot count equals the Betti polynomial
(10). The primitive lattice and rank bound give, degree by degree,
\[
 Q_t(m)\subseteq P_t(m)\subseteq R_{A,\ell}(m),\qquad
 \rank Q_t(m)=\rank P_t(m)=\rank H^{2m}(W_t;\Z).
\]
A primitive sublattice contains every integral point of its rational
span. Hence \(P_t(m)=Q_t(m)\): all prefixes are primitive.
The adapted-basis lemma after (8) now proves (4) in every rank.

\paragraph{A consequence, not a circular hypothesis.}
Restriction \(H^*(\Gr(\ell,n);\Z)\to H^*(W_t;\Z)\) is surjective.
Indeed the incidence pullback is primitive, and the composite has
primitive image of the full Betti rank.
Transposition preserves Smith invariants, so the primitive coefficient
row lattice gives the required primitive image as well.
Chern generation is obtained after the lattice comparison.

\paragraph{Checks and audit outcome.}
The fresh review included the earlier prefix reduction, both primitive
pullbacks, function-field cycle lifts, bundle truncations, energy minor,
both parity identities, and endpoint cases. No gap was found.
Existing exact records check 585,243 corner/bead comparisons in 72
parameter cases and 21 fresh integer LR module comparisons through
\(\ell=8\); independent graph proofs also cover the three highest
layers. A fresh pure-Python checker independently passed 16,200,912
corner comparisons on 1,430,949 partitions in nine larger cases,
\(n=14,\ldots,22\), \(\ell=7,\ldots,11\), with 88 exact weighted
polynomial comparisons. Earlier records check 174 Smith lists,
430 prefix ranks, and 516 integer row-module equalities.
These corroborate transcription and implementation; the unrestricted
proof is the argument above.
No external expert review or formal proof certificate is claimed.

\newpage
\section{Published formulas and the remaining type-D problem}
\published The polynomial \(U\) in (3) is classical Birkhoff/Hall
enumeration \cite{Birkhoff,Sage}, also counting invariant subspaces of
a maximal-rank square-zero map.
More general operator-profile counts are already covered by Ram
\cite{Ram}. For \(N^2=0\), defect \(e\) is profile \((\ell,e)\).
Counting those objects alone does not prove the Smith assertion.

There is also a published complex representation-theoretic match.
Chari--Schneider--Shereen--Wand \cite{CSSW},
Theorem 3.3(ii)--(iii), give
\[
 [W_{\rm loc}(n\omega):D(2,(n-2e)\omega)]_q=q^{ec}G(r,e).
\]
Their character identity gives \(U_{n-2e,\ell-e}\) for the weight
\(n-2\ell\) subspace of the remaining Demazure module.
Thus the expression in (4) already has a character interpretation.
That work is over \(\C\); identifying this polynomial with the
integral Springer-image layers requires the additional comparison in
Sections 4--8. We make no originality claim for that identification.

\paragraph{Relation to the Euler generator in Part III.}
\deduction The Euler sector generated by \(p f(q)\) on
\(\OG(a,2n)\), for \(a\ge2\), has the Smith factors of the normalized B
image on \(\OG(a-1,2n-1)\), shifted by polynomial degree \(k\).
This uses the explicit integral difference basis in BKT
Proposition 2 \cite{BKTD}. Formula (4) now determines those sector
multiplicities in every rank.

The full D image needs further gluing data. A swapped pair of typed
Schubert classes has sum and difference of index two.
The companion two-flag lattice lemma gives, for a paired
sector-exponent pair \((i,j)\), exponents
\[
 \bigl(\min(i,j),\ \max(i,j)+1\bigr).
 \tag{19}
\]
Relative flags over \(\F_2\) determine the pairing.
The two sector Smith lists alone do not determine it.
\conjectural The general D square-sector equality with B remains
unproved here. A complete type-D partition classification remains open
in these notes; the new B argument does not settle that gluing problem.

\paragraph{Literal character forms.}
The chosen normalized forms are those used in the proofs of Parts II
and III. For the defining/vector SO/Sp character forms in Part I
\cite{K}, a degree-\(d\) image is \(2^d\) times its normalized
counterpart. Every nonzero factor is therefore multiplied by \(2^d\);
in B degree \(2m\), \(2^e\) becomes \(2^{2m+e}\).
This statement does not automatically extend to a Spin character lattice.

The complete general proof, fresh dependency audit, and exact reports
are linked in \cite{Repo}. The unresolved D statements and literature
limits remain labeled separately from the proved B deduction.

\newpage
\renewcommand{\refname}{References and companion material}
{\small
\setlength{\parskip}{2pt}
\begin{thebibliography}{BKT-IG}
\bibitem[KR20]{KR}
S. Kumar and S. Rogers,
\emph{Representation ring of Levi subgroups versus cohomology ring of flag
varieties II}, J. Algebra \textbf{556} (2020), 340--362.
Theorems 16 and 29; integral form in the proof of Theorem 16, p.~353.
\href{https://kumar.math.unc.edu/papers/kumar81.pdf}{Author PDF};
\href{https://doi.org/10.1016/j.jalgebra.2020.02.029}{DOI}.

\bibitem[KX24]{KX}
S. Kumar and J. Xie,
\emph{Representation ring of Levi subgroups versus cohomology ring of flag
varieties III}, J. Pure Appl. Algebra \textbf{228} (2024), 107539.
Theorem 3.6, Proposition 3.7, Theorem 3.8, equations (28)--(33).
\href{https://kumar.math.unc.edu/papers/kumar86.pdf}{Author PDF}.

\bibitem[K16]{K}
S. Kumar,
\emph{Representation ring of Levi subgroups versus cohomology ring of flag
varieties}, Math. Ann. \textbf{366} (2016), 395--415.
Proposition 12(b), Lemma 15(a), equation (15).
\href{https://kumar.math.unc.edu/papers/kumar76.pdf}{Author PDF}.

\bibitem[BKT-IG]{BKTIG}
A. S. Buch, A. Kresch and H. Tamvakis,
\emph{A Giambelli formula for isotropic Grassmannians}.
Definition 1, Theorem 1, Sections 2.1 and 6.1.
\href{https://math.umd.edu/~harryt/papers/isogiamfinal.pdf}{Author text}.

\bibitem[BKT-D]{BKTD}
A. S. Buch, A. Kresch and H. Tamvakis,
\emph{A Giambelli formula for even orthogonal Grassmannians}.
Proposition 2, Theorem 4, Section 5.2.
\href{https://math.umd.edu/~harryt/papers/oggiam.pdf}{Author text}.

\bibitem[PR]{PR}
P. Pragacz and J. Ratajski,
\emph{Formulas for Lagrangian and orthogonal degeneracy loci;
\(\widetilde Q\)-polynomial approach}.
Appendix A, Sections 10.3 and 10.5.
\href{https://www.impan.pl/~pragacz/download/comp3868.pdf}{Author text}.

\bibitem[PW]{PW}
I. Panin and C. Walter, \emph{Quaternionic Grassmannians and Borel classes}.
Theorems 11.2--11.3 (relative basis and presentation).
\href{https://arxiv.org/pdf/1011.0649}{Author text}.

\bibitem[W]{Wildon}
M. Wildon, \emph{A generalized SXP rule proved by bijections and
involutions}. Equation (1), the Littlewood/SXP identity.
\href{https://arxiv.org/pdf/1508.07030}{Author text}.

\bibitem[CSSW14]{CSSW}
V. Chari, L. Schneider, P. Shereen and J. Wand,
\emph{Modules with Demazure Flags and Character Formulae},
SIGMA \textbf{10} (2014), 032.
Section 2.1; Theorem 3.3(ii)--(iii); Remark 3.4.
\href{https://arxiv.org/pdf/1310.5191}{Published author text};
\href{https://doi.org/10.3842/SIGMA.2014.032}{DOI}.

\bibitem[B35]{Birkhoff}
G. Birkhoff, \emph{Subgroups of Abelian Groups},
Proc. London Math. Soc. \textbf{38} (1935), 385--401.
\href{https://doi.org/10.1112/plms/s2-38.1.385}{Bibliographic record}.
The original full text was not accessed in the repository audit; the
general formula was inspected in \cite{Sage} and its specialization
checked against both Sage counting algorithms.

\bibitem[Sage]{Sage}
SageMath reference, \emph{\(q\)-analogues},
\texttt{q\_subgroups\_of\_abelian\_group}.
The Birkhoff formula is displayed explicitly, citing Butler (1987).
\href{https://doc.sagemath.org/html/en/reference/combinat/sage/combinat/q_analogues.html}{Official documentation}.

\bibitem[R25]{Ram}
S. Ram, \emph{Subspace profiles, \(q\)-Whittaker functions and Krylov methods},
Sem. Lothar. Combin. \textbf{93B} (2025), article 115.
Definition 1.1 and Theorem 5.2 (also Theorem 1.3).
\href{https://www.mat.univie.ac.at/~slc/wpapers/FPSAC2025/115.pdf}{Conference paper}.

\bibitem[Repo]{Repo}
\emph{Sage-programs: general type-B proof and integral image notes}.
Unpublished companion material, audited 19 October 2026.
\href{https://github.com/sean-rogers-t/Sage-programs/blob/5e718334a0f8ed9b31166aa3bf42185a60d05c34/borel/TYPE_B_GENERAL_SMITH.md}{Full general proof at audited baseline \texttt{5e71833}};
\href{https://github.com/sean-rogers-t/Sage-programs/blob/codex/borel-library/borel/TYPE_B_GENERAL_AUDIT.md}{fresh dependency-by-dependency audit};
\href{https://github.com/sean-rogers-t/Sage-programs/blob/codex/borel-library/borel/PROOF_STATUS.md}{current claim and source status}.
The linked notes retain the full prefix, normalization, type-D, and
independent verification arguments. The baseline link is immutable;
the two branch links may be updated.
\end{thebibliography}
}
\end{document}
